\documentclass[11pt]{article}

\usepackage[margin=1in]{geometry}
\usepackage{amsmath,amssymb,amsthm}
\usepackage{mathtools}
\usepackage{enumitem}
\usepackage{hyperref}
\usepackage{natbib}

\newtheorem{proposition}{Proposition}
\newtheorem{corollary}{Corollary}
\newtheorem{definition}{Definition}
\newtheorem{assumption}{Assumption}
\newtheorem{remark}{Remark}
\newtheorem{lemma}{Lemma}

\title{Strategic Under-Specification: A Stackelberg Game Between User and Assistant}
\author{Pranjal Agarwal \\ \small BITS Pilani, Hyderabad Campus}
\date{}

\begin{document}
\maketitle

\begin{abstract}
Users of coding assistants routinely under-specify their intent, and the assistant's
policy for handling missing information---guess silently, or ask a clarifying
question---is itself a strategic choice that users respond to. We formalize this as
a Stackelberg game: the assistant commits to an ask/guess policy, and a population
of users with privately known, heterogeneous specification costs best-respond. We
first solve a decision-theoretic benchmark in which the policy is exogenous, and
show that improving the assistant's raw guessing accuracy alone cannot make users
worse off, even though it induces them to specify less---so naive ``better guessing
causes worse specs and worse outcomes'' stories are not generically true. We then
close the loop and solve for the assistant's own optimal policy under two regimes:
a \emph{pooling} policy restricted to a single population-wide ask rate, and a
\emph{screening} policy that can condition on the user's declared specification. We
show that an unbiased (welfare-aligned) assistant achieves full efficiency under
screening despite users' private information---incentive compatibility is free when
principal and agent share an objective---so all screening-regime inefficiency is
attributable to misalignment between the assistant's effective objective and true
user welfare, not to private information per se. Pooling, by contrast, loses
efficiency from restricted instruments even when unbiased. We characterize both
equilibria in closed or near-closed form under a tractable linear-risk
specification, derive testable comparative statics, and lay out a simulation and
human-study design to test them.
\end{abstract}

\section{Introduction}

A user typing a coding request faces a choice: spend effort articulating edge
cases, naming conventions, and constraints, or leave them implicit and trust the
assistant to infer them. The assistant, in turn, must decide for each unspecified
detail whether to guess or to ask. Neither choice is made in a vacuum: the user's
specification effort depends on what they expect the assistant to do with silence,
and the assistant's ask/guess policy is chosen (explicitly, or as an emergent
property of training) with some implicit model of how users will respond. Section 2
treats the assistant's policy as exogenous and solves the user's problem against
it---a one-sided optimization, useful as a benchmark but not yet a game. Sections
4--7 close the loop: the assistant is a Stackelberg leader that commits to a policy
anticipating the population's best response, and we solve for that policy under two
different instrument restrictions, then compare them.

\section{Environment and the Decision-Theoretic Benchmark}
\label{sec:baseline}

\subsection{Setup}

A task has $k$ binary, independent, intent-relevant attributes
$\theta = (\theta_1,\dots,\theta_k) \in \{0,1\}^k$, each drawn i.i.d.\ uniform and
privately known to the user. The user chooses a specification effort level
$m \in \{0,\dots,k\}$ and truthfully reveals $m$ attributes.\footnote{We assume
truthful revelation conditional on the choice to specify: the strategic margin is
\emph{how much} to say, not whether to lie about what is said. A cheap-talk layer
on top of this costly-specification layer is a natural extension; see
Section~\ref{sec:discussion}.} Specifying $m$ attributes costs $c(m)$, with
$c(0)=0$, $c' > 0$, $c'' \ge 0$.

For each of the $k-m$ unspecified attributes, the assistant either \emph{asks}
(revealing the true value, at a cost to the user) or \emph{guesses}, succeeding
independently with probability $g \in (0,1)$. We model the assistant's policy at
this stage as a uniform ask rate $a \in [0,1]$ applied to every unspecified
attribute, so each is correctly resolved with probability
$q(a,g) = a + (1-a)g$. Task success requires \emph{all} $k$ attributes to be
correctly resolved (conjunctive: one wrong parameter or convention typically breaks
the deliverable). The user's payoff is
\[
U(m; a, g) \;=\; V \, q(a,g)^{\,k-m} \;-\; c(m) \;-\; a(k-m)\kappa ,
\]
with $V>0$ the value of success and $\kappa>0$ the user's cost of answering one
clarifying question. We treat $m$ as continuous on $[0,k]$ for tractability.

\subsection{Proposition 1: accuracy alone does not create moral hazard}

\begin{proposition}\label{prop:accuracy}
Fix $a$. Let $U^*(g) = \max_{m \in [0,k]} U(m;a,g)$. Then $U^*$ is weakly
increasing in $g$.
\end{proposition}

\begin{proof}
By the envelope theorem, at the optimum $m^*(g)$,
$dU^*/dg = \partial U/\partial g\big|_{m=m^*(g)} = V(k-m^*)q^{k-m^*-1}(1-a) \ge 0$.
The envelope theorem lets us ignore $dm^*/dg$: because $m^*$ is chosen optimally,
its first-order response to $g$ has no first-order effect on $U^*$.
\end{proof}

\begin{assumption}[Regularity]\label{ass:reg}
In the relevant range of $m$, $(k-m)\ln q(a,g) \le -1$.
\end{assumption}

\begin{corollary}\label{cor:mono}
Under Assumption~\ref{ass:reg}, $m^*(g)$ is weakly decreasing in $g$ (decreasing
differences in $(m,g)$; Topkis's theorem).
\end{corollary}

So, holding the assistant's policy fixed, users specify less as the assistant gets
better at guessing, yet are never worse off for it. This is the correct benchmark
against which any claimed moral hazard should be measured, and it isolates the
question this paper is really about: what happens once the policy $a$ is no longer
exogenous, but chosen by a strategic assistant that anticipates this response?

\section{Sources of Strategic Under-Specification}

Proposition~\ref{prop:accuracy} assumed the assistant's ask rate is fixed and that
users correctly anticipate $q(a,g)$. Both assumptions can fail.

\subsection{Cost externality}
If a third party (a reviewer, downstream production systems) bears part of the
cost of a wrong guess, the user's private $(\kappa, V)$ understate the social
values, and the user under-specifies relative to the social optimum at any fixed
$g$---worsening, by Corollary~\ref{cor:mono}, as $g$ rises. Standard; not our focus.

\subsection{Miscalibrated belief updating under silent failure}
If wrong guesses are not always detected, and detection probability $d(m)$ is
increasing in $m$, users infer $g$ from experience rather than observing it, and
under-specifying users are also less likely to catch their own errors---an
upward-biased $\hat g$ that reinforces under-specification. This is a dynamic
learning story rather than a static equilibrium result; we return to it in the
empirical design (prediction P5).

\subsection{Belief about the assistant's policy}
\label{sec:naive-soph}
Let users be \emph{sophisticated} (correctly infer the assistant's true ask rate
$a_A$) or \emph{naive} (believe the assistant follows the welfare-optimal rate
$a^\dagger \ge a_A$).

\begin{proposition}\label{prop:bias}
Under Assumption~\ref{ass:reg}, naive users choose
$m^*_{\text{naive}} \le m^*_{\text{soph}}$ and
$U(m^*_{\text{naive}};a_A,g) < U(m^*_{\text{soph}};a_A,g)$: naive users
under-specify and are strictly worse off than sophisticated users, holding raw
guessing accuracy $g$ fixed throughout.
\end{proposition}

\begin{proof}
Corollary~\ref{cor:mono} gives $m^*(\cdot)$ non-increasing in its accuracy
argument; $m^*_{\text{soph}}$ is optimal against the true $q(a_A,g)$, so any other
$m$, including $m^*_{\text{naive}}$, delivers strictly lower $U(\cdot;a_A,g)$
whenever it differs.
\end{proof}

This is the source of genuine moral hazard identified so far: not raw accuracy, but
an \emph{unobserved gap between the user's belief about the assistant's policy and
the assistant's true policy}. Section~\ref{sec:model1} solves for what that true
policy $a_A$ actually is when the assistant is itself a strategic, possibly biased,
Stackelberg leader---so Proposition~\ref{prop:bias} can be instantiated with a
derived $a_A$ rather than an assumed one.

\section{The Assistant's Policy as a Stackelberg Game}
\label{sec:stackelberg}

\subsection{General formulation}

The assistant (leader) commits ex ante to a policy $\pi$ before meeting any user.
A population of users (followers) has privately known specification-cost type
$\kappa \sim F$ on $[\kappa_L,\kappa_H]$, with cost $c(m;\kappa)=\kappa m^2/2$, and
best-responds $m^*(\kappa;\pi) = \arg\max_m U(m;\pi,\kappa)$. The leader's own
objective $\Pi$ need not equal $U$: we write
\[
\Pi_\kappa(m,a) \;=\; \mu_A\, U(m,a;\kappa) \;-\; (\lambda_A - c_Q)\, a\,(k-m),
\]
where $\mu_A \in (0,1]$ discounts the leader's weight on true (downstream) user
welfare, $c_Q$ is the user's true cost of answering one clarifying question, and
$\lambda_A$ is the leader's own \emph{perceived} per-question friction cost, which
may differ from $c_Q$ (e.g., a training signal that penalizes extra turns more
than users actually mind answering a direct question). The leader is
\emph{unbiased} iff $\mu_A=1$ and $\lambda_A=c_Q$, in which case $\Pi_\kappa \equiv
U(\cdot;\kappa)$ exactly. The leader solves
$\pi^* = \arg\max_\pi \mathbb{E}_{\kappa\sim F}[\Pi_\kappa(m^*(\kappa;\pi),\pi)]$,
a bilevel program.

\begin{remark}[Existence]
For Model I below ($\pi=a\in[0,1]$, a compact scalar policy), $m^*(\kappa;a)$ is
continuous in $a$ (strict concavity of $U$ in $m$ plus the implicit function
theorem), so $\mathbb{E}_\kappa[\Pi_\kappa]$ is continuous in $a$ and a maximizer
exists on $[0,1]$ by Weierstrass's theorem. For Model II below (a finite menu over
two types), the feasible set is a compact subset of $\mathbb{R}^4$ cut out by
linear IC/IR constraints and a continuous objective, so a maximizer again exists
by Weierstrass's theorem; no ironing or ex-ante randomization is needed with only
two types.
\end{remark}

For tractability in what follows we replace the exact conjunctive success
probability $q^{k-m}$ with its first-order (small-risk) approximation around
$q\to 1$: expected loss $\approx L(1-q)(k-m)$, where $L$ is the loss from a single
unresolved-and-incorrect attribute. (This linear-risk form is already implicit in
the stakes notation $L(s)$ used informally in earlier drafts of this model; here we
make it explicit and load-bearing.) With $\Lambda \equiv L(1-g)$, the user's payoff
becomes
\[
U(m,a;\kappa) \;=\; V \;-\; \Lambda(1-a)(k-m) \;-\; \frac{\kappa}{2}m^2 \;-\; a\,c_Q\,(k-m).
\]

\subsection{The first-best benchmark}
\label{sec:firstbest}

Before analyzing either regime, characterize the fully-informed, fully-aligned
optimum: $(m^{FB}(\kappa), a^{FB}(\kappa)) = \arg\max_{m,a} U(m,a;\kappa)$.

\begin{proposition}[First-best is a threshold rule in $\kappa$]\label{prop:firstbest}
$U$ is linear in $a$ for fixed $m$, so $a^{FB}(\kappa) \in \{0,1\}$. Define
$\kappa^* = (\Lambda+c_Q)/(2k)$. If $\Lambda > c_Q$, then $a^{FB}(\kappa)=1$
(always ask) for $\kappa > \kappa^*$ and $a^{FB}(\kappa)=0$ (never ask, guess) for
$\kappa < \kappa^*$.
\end{proposition}

\begin{proof}
At $a=1$: $m^*(\kappa;1)=c_Q/\kappa$ and $U(m^*,1;\kappa)=V-c_Qk+c_Q^2/(2\kappa)$.
At $a=0$: $m^*(\kappa;0)=\Lambda/\kappa$ and
$U(m^*,0;\kappa)=V-\Lambda k+\Lambda^2/(2\kappa)$. Subtracting,
\[
U(a{=}1)-U(a{=}0) \;=\; [\Lambda-c_Q]\left(k - \frac{\Lambda+c_Q}{2\kappa}\right),
\]
which is positive iff $\Lambda>c_Q$ and $\kappa>(\Lambda+c_Q)/(2k)=\kappa^*$.
\end{proof}

This is a clean, closed-form, type-dependent first best: high-cost (terse) users,
who cannot cheaply self-provide detail, should always be asked; low-cost (verbose)
users, who can cheaply write enough that the small residual is safe to guess,
should never be asked. This threshold rule replaces and generalizes the
stakes-weighted threshold from earlier drafts of this model (that result is the
special case obtained by varying $L(s)$ at fixed $\kappa$ instead of varying
$\kappa$ at fixed $L$).

\section{Model I: A Pooling (Scalar) Policy}
\label{sec:model1}

The leader is restricted to a single population-wide ask rate $a\in[0,1]$---it
cannot condition on anything user-specific. This is the natural benchmark for a
policy that is fixed at training time and cannot observe individual users.

\subsection{User best response}

The FOC $\partial U/\partial m = \Lambda(1-a) - \kappa m + ac_Q = 0$ gives, with
$\beta(\kappa)=\Lambda/\kappa$ and $\gamma(\kappa)=(\Lambda-c_Q)/\kappa$,
\[
m^*(\kappa;a) \;=\; \beta(\kappa) - a\,\gamma(\kappa), \qquad
\frac{\partial m^*}{\partial a} \;=\; -\gamma(\kappa) \;=\; \frac{c_Q-\Lambda}{\kappa}.
\]

\begin{remark}[Non-monotone effect of ask rate on specification]
Unlike Corollary~\ref{cor:mono}, the sign of $\partial m^*/\partial a$ here does
not depend on type: it is positive for \emph{every} type iff $c_Q > \Lambda$, i.e.,
iff answering a question costs users more than the expected loss it avoids. When
that holds, a higher population ask rate induces users to specify \emph{more}, not
less---they pre-empt costly clarifying questions by writing the answer down
themselves. This is a genuinely new force relative to Section~\ref{sec:baseline},
where $a$ only entered through accuracy, and is a distinct, directly testable
prediction (P1$'$ below).
\end{remark}

\subsection{The leader's problem}

Let $R(a,\kappa)=k-m^*(\kappa;a)$. The leader's per-type payoff is
$\Pi_\kappa(a) = \mu_A V - C(a)R(a,\kappa)$ with $C(a)=\mu_A\Lambda + a(\lambda_A -
\mu_A\Lambda)$. Averaging over $F$ with $\bar R(a) = \bar R_0 + a\bar\gamma$
(where $\bar R_0=k-\mathbb{E}[\beta]$, $\bar\gamma=\mathbb{E}[\gamma]$), $\Pi(a)$
is quadratic in $a$:
\[
\Pi(a) \;=\; \mu_A V - C_0\bar R_0 - a\big[C_0\bar\gamma + \Delta\bar R_0\big] -
a^2\Delta\bar\gamma, \qquad C_0=\mu_A\Lambda,\;\; \Delta=\lambda_A-\mu_A\Lambda.
\]

\begin{proposition}[Stackelberg-optimal pooling rate]\label{prop:aSE}
Provided $\Delta\bar\gamma>0$ (second-order condition), the leader's optimal
pooling ask rate is
\[
a^{SE} \;=\; -\,\frac{C_0\bar\gamma + \Delta \bar R_0}{2\Delta\bar\gamma}.
\]
\end{proposition}

\begin{corollary}[Bias shifts the pooling rate]
$\partial a^{SE}/\partial \lambda_A$ has the sign of $2\bar\gamma a^{SE} - \bar
R_0$ (by implicit differentiation of the FOC). Under a further regularity
condition analogous to Assumption~\ref{ass:reg} (the residual $\bar R(a^{SE})$
stays above half of its $a{=}0$ level), $a^{SE}$ is strictly decreasing in
$\lambda_A$: an assistant whose training penalizes clarifying questions more than
users actually mind them ($\lambda_A>c_Q$) sets a strictly lower pooling ask rate
than an unbiased assistant would.
\end{corollary}

\begin{corollary}[Pooling is a corner solution, not a compromise]\label{cor:corner}
In the unbiased case ($\mu_A=1,\lambda_A=c_Q$), $\Delta=c_Q-\Lambda$ and
$\bar\gamma=(\Lambda-c_Q)\mathbb{E}[1/\kappa]$, so
\[
\Delta\bar\gamma \;=\; -(\Lambda-c_Q)^2\,\mathbb{E}[1/\kappa] \;\le\; 0 \quad
\text{always}.
\]
The second-order condition of Proposition~\ref{prop:aSE} therefore \emph{fails}
in the unbiased case: $\Pi(a)$ is weakly convex on $[0,1]$, its stationary point
is a minimum, and the true optimum is a corner, $a^{SE}\in\{0,1\}$, not an
interior compromise as an earlier draft of this result claimed. Comparing the
two corners directly,
\[
\Pi(1)-\Pi(0) \;=\; (\Lambda-c_Q)\Big[\bar R_0 - c_Q\,\mathbb{E}[1/\kappa]\Big],
\]
so the unbiased pooling optimum is $a^{SE}=1$ (always ask) when this is positive
and $a^{SE}=0$ (never ask) when negative. Pooling therefore does not average
across the population's heterogeneous first-best preferences
(Proposition~\ref{prop:firstbest}); it serves one type's preferred extreme
exactly and gives the other type the extreme it least wants---a starker
inefficiency than a blended compromise, and one that sharpens rather than
weakens the contrast with screening in Section~\ref{sec:compare}: screening's
value is not merely ``less averaging'' but the ability to treat types
differently at all. (An interior $a^{SE}$ remains possible away from the
unbiased case, specifically when $\lambda_A-\mu_A\Lambda$ and $\Lambda-c_Q$
share a sign, so bias can change not only the level of $a^{SE}$ but whether
pooling compromises or picks a winner.)
\end{corollary}

This last point matters for the comparison in Section~\ref{sec:compare}: Model I
has two independent, additive sources of loss relative to first best---the
pooling restriction, and bias.

\section{Model II: A Screening (Menu) Policy}
\label{sec:model2}

Now let the leader condition its ask rate on the user's declared specification
level (equivalently, by the revelation principle, offer a menu of bundles that
users self-select into). Consider two types $\kappa_L<\kappa_H$ with population
shares $f_L,f_H$, and let the leader choose bundles $(m_L,a_L),(m_H,a_H)$ subject
to
\[
\text{IC}_\kappa:\;\; U(m_\kappa,a_\kappa;\kappa) \ge U(m_{\kappa'},a_{\kappa'};\kappa), \qquad
\text{IR}_\kappa:\;\; U(m_\kappa,a_\kappa;\kappa) \ge \underline{U},
\]
to maximize $f_L\Pi_{\kappa_L}(m_L,a_L)+f_H\Pi_{\kappa_H}(m_H,a_H)$.

\subsection{Single-crossing}
The marginal rate of substitution between $m$ and $a$ along an indifference curve
is $\text{MRS}(\kappa) = -\partial U/\partial m \big/ \partial U/\partial a$; using
the FOC expressions above, $\partial \text{MRS}/\partial \kappa$ has a consistent
sign across the relevant range (higher-$\kappa$ types have steeper indifference
curves in $(m,a)$-space: they require more compensating ask-service per unit of
specification given up). This is the standard Spence--Mirrlees condition and
underwrites everything below; we omit the routine but lengthy verification.

\subsection{No distortion without bias}

\begin{proposition}[Efficiency survives private information when the leader is unbiased]
\label{prop:nodistortion}
If $\mu_A=1$ and $\lambda_A=c_Q$ (so $\Pi_\kappa \equiv U(\cdot;\kappa)$ for every
$\kappa$), then implementing each type's own first-best bundle,
$(m_\kappa,a_\kappa)=(m^{FB}(\kappa),a^{FB}(\kappa))$ from
Proposition~\ref{prop:firstbest}, is incentive compatible, and no further
mechanism design can do better for the leader.
\end{proposition}

\begin{proof}
$(m^{FB}(\kappa),a^{FB}(\kappa))$ maximizes $U(\cdot;\kappa)$ by definition, so for
any $\kappa'\ne\kappa$, $U(m^{FB}(\kappa),a^{FB}(\kappa);\kappa) \ge
U(m^{FB}(\kappa'),a^{FB}(\kappa');\kappa)$ trivially: no bundle, including another
type's, can give type $\kappa$ more utility than its own maximizer. IC holds with
no distortion needed. Since $\Pi_\kappa\equiv U(\cdot;\kappa)$, the leader's payoff
from serving each true type is exactly maximized at that type's own first best, so
this menu is also leader-optimal.
\end{proof}

This is the paper's sharpest result: unlike classical regulation (Baron--Myerson,
Laffont--Tirole), where the principal pays a literal transfer and therefore has an
independent reason to dislike conceding rent, here the leader's payoff can
coincide exactly with user welfare, and when it does, incentive compatibility is
\emph{free}---private information costs nothing. Distortion in this setting is
therefore a pure signature of misalignment, not of asymmetric information per se.

\subsection{Distortion under bias}

When $\Pi_\kappa \ne U(\cdot;\kappa)$ (bias present), the leader's preferred bundle
per type, $(m_\kappa^{B},a_\kappa^{B})=\arg\max_{m,a}\Pi_\kappa(m,a)$ (computable
in closed form by the same method as Proposition~\ref{prop:firstbest}, with
$\Lambda,c_Q$ replaced by their $\mu_A,\lambda_A$-weighted counterparts), no longer
automatically satisfies IC, because $\Pi_\kappa$ and the user's true $U(\cdot;\kappa)$
now rank bundles differently.

\begin{proposition}[Downward distortion of the high-cost type's service]
\label{prop:distortion}
Under the single-crossing condition and the empirically motivated under-asking
bias ($\lambda_A > \mu_A c_Q$, i.e.\ the leader over-weights the friction of
asking relative to true user cost), the leader's optimal incentive-compatible
menu sets $a_L^{SB}=a_L^{B}$ (no further distortion for the low-cost type, whose
biased first-best is already the binding corner $a=0$) and $a_H^{SB} < a_H^{B} \le
a_H^{FB}=1$ (the high-cost type is asked strictly less than even the leader's own
already-biased preference for them), with $a_H^{SB}$ characterized by the
standard virtual-surplus first-order condition
$\partial\Pi_{\kappa_H}/\partial a_H = \frac{f_L}{f_H}\big[\kappa_H-\kappa_L\big]
m_H \cdot \partial m_H/\partial a_H$ (rent-adjustment term on the right).
\end{proposition}

\begin{proof}[Proof sketch, corrected]
Relax the problem to IC$_L$ alone binding. Unlike classical transfer-based
screening (Baron--Myerson, Laffont--Tirole), where the principal's payoff is
directly reduced by the rent conceded to the efficient type---so the standard
technique pushes the inefficient type's IR down to reservation to minimize that
rent---here $\Pi_{\kappa_H}$ already contains $\mu_A U(\cdot;\kappa_H)$ directly,
so the leader's own payoff from serving the high-cost type \emph{rises} with
that type's true welfare: there is no independent budget motive to depress
$U_H$ to $\underline{U}$, and IR$_H$ need not bind. The bundle given to the
high-cost type is instead set by its own (biased) unconstrained optimum
$(m_H^{B},a_H^{B})=\arg\max\Pi_{\kappa_H}$, moved only as far as needed to
satisfy IC$_L$---i.e., only until $U(m_H,a_H;\kappa_L)=U(m_L,a_L;\kappa_L)$
binds---using the envelope representation of the low type's temptation to
mimic, $U(m_H,a_H;\kappa_L)-U(m_H,a_H;\kappa_H)=(\kappa_H-\kappa_L)m_H^2/2$.
This still yields $a_H^{SB}<a_H^{B}$ whenever the unconstrained biased optimum
would otherwise violate IC$_L$, but the resulting first-order condition is not
the classical virtual-surplus expression a naive import of Baron--Myerson would
suggest; see the remark below.
\end{proof}

\begin{remark}[Correction from numerical audit]
An earlier draft of this proof sketch assumed, by direct analogy with
Baron--Myerson, that IC$_L$ and IR$_H$ bind while IC$_H$ and IR$_L$ are slack.
Solving the full four-variable constrained problem numerically instead finds
IC$_L$ binding alone, with IR$_H$, IC$_H$, and IR$_L$ slack at representative
under-asking-bias parameters, for exactly the reason given in the proof sketch
above---and it is the same class of error already corrected once in this model
(Proposition~\ref{prop:nodistortion}): a classical mechanism-design result that
relies on the principal's payoff being reduced by conceded rent does not
transfer mechanically to a principal whose payoff already contains the agent's
utility. We flag the exact active constraint set as parameter-region-specific
rather than proven for all bias magnitudes---for sufficiently extreme bias,
$\Pi_{\kappa_H}$ and $U(\cdot;\kappa_H)$ could diverge enough that IC$_H$ binds
instead, which would change the character of the result. The claim we retain
with confidence is qualitative: some constraint binds and produces
$a_H^{SB}<a_H^{B}$ under under-asking bias; the exact active set should be
checked numerically for the parameter region of interest rather than assumed
from the classical template.
\end{remark}

\begin{remark}
The direction in Proposition~\ref{prop:distortion} is specific to the
under-asking bias case ($\lambda_A>\mu_A c_Q$), which we take to be the
empirically relevant one (Section~\ref{sec:naive-soph}'s motivating story: turn-level
satisfaction signals penalize extra clarifying questions more than users actually
mind answering them). Under the opposite bias, the direction of distortion can
flip; we do not claim a sign that holds for all $\Pi_\kappa$.
\end{remark}

\section{Comparing the Two Regimes}
\label{sec:compare}

\begin{corollary}[Menus weakly dominate pooling]
The leader's optimal expected payoff under Model II is weakly greater than under
Model I: any pooling policy $a$ is a feasible (degenerate) menu with
$a_L=a_H=a$, trivially satisfying IC, so the unconstrained menu optimum is at
least as good.
\end{corollary}

The more interesting comparison is \emph{where the two regimes' losses come from},
which Propositions~\ref{prop:aSE}--\ref{prop:distortion} let us decompose cleanly:

\begin{itemize}[leftmargin=*]
\item \textbf{Model I (pooling)} loses welfare from two sources that need not be
additive in the same simple sense as in earlier drafts: (i) the pooling
restriction itself, present even when $\mu_A{=}1,\lambda_A{=}c_Q$, where it
manifests as a corner solution that fully serves one type and fully fails the
other (Corollary~\ref{cor:corner}), and (ii) bias, which shifts $a^{SE}$
(Corollary 2) and can additionally flip pooling from a corner into an interior
compromise or vice versa, depending on the sign alignment noted in
Corollary~\ref{cor:corner}. The two sources interact rather than simply add.
\item \textbf{Model II (screening)} loses welfare from bias \emph{only}
(Proposition~\ref{prop:nodistortion}): with no bias, screening recovers full
efficiency despite private information. With bias, the loss is the
rent-extraction-driven downward distortion of Proposition~\ref{prop:distortion},
which is a genuinely different mechanism from Model I's pooling loss (it doesn't
arise from being unable to condition on anything---it arises from not being able
to trust what is conditioned on).
\end{itemize}

\begin{remark}[Compounding, testable]
Because Model II's distortion is triggered by, and scales with, the same bias
parameters ($\mu_A,\lambda_A$ vs.\ $c_Q$) that drive Model I's shift in
$a^{SE}$, reducing measured assistant bias should improve outcomes for high-cost
(terse) users under \emph{both} regimes, but the welfare \emph{gap between}
menu-conditioned and pooled policies should track population heterogeneity in
$\kappa$ (the spread $\kappa_H-\kappa_L$ and how far $F$ sits on either side of
$\kappa^*$) almost independently of bias. This gives two separable, testable
quantities from the same field data: a bias-sensitive term and a
heterogeneity-sensitive term.
\end{remark}

\section{Testable Predictions}
\label{sec:predictions}

\begin{enumerate}[leftmargin=*]
\item[\textbf{P1}] (Prop.~\ref{prop:accuracy}, Cor.~\ref{cor:mono}) Holding the
assistant's policy fixed, higher guessing accuracy should be associated with
shorter specifications but higher or equal success rates---never both shorter
specs and lower success.
\item[\textbf{P1$'$}] (Model I remark) Holding accuracy fixed, the effect of a
higher \emph{ask rate} on specification length should flip sign depending on
whether $c_Q \gtrless \Lambda=L(1-g)$: where clarifying questions are cheap
relative to the risk they resolve, users should specify less as ask rate rises
(substitution); where questions are relatively costly to answer, users should
specify \emph{more} (pre-emption). This is a within-model prediction distinct
from P1's accuracy channel.
\item[\textbf{P2}] (Prop.~\ref{prop:bias}) The correlation between turn-level
satisfaction and eventual task success should weaken, or invert, specifically for
under-asking assistants, and can now be benchmarked against the derived $a^{SE}$
of Proposition~\ref{prop:aSE} rather than an assumed ask rate.
\item[\textbf{P3}] Making ask-rate statistics visible to users should shift
specification length toward the sophisticated best response of
Section~\ref{sec:naive-soph}.
\item[\textbf{P4}] (Prop.~\ref{prop:firstbest}) Ask rates should rise with a
user's revealed difficulty specifying detail (terseness), sharply around a
threshold $\kappa^*$, rather than responding only to model confidence.
\item[\textbf{P5}] Users with lower self-reported error-catching rates should show
declining specification effort over a session (Section 3.2's dynamic channel).
\item[\textbf{P6}] (Cor.~3) Under a pooling policy, terse users should be
under-served and verbose users over-served relative to \emph{each type's own}
first best, even for a demonstrably unbiased assistant---a gap attributable
purely to instrument restriction, distinguishable in data from bias by holding
$\lambda_A$ fixed (e.g., via an A/B test on friction penalties in training) and
varying only whether the policy can condition on declared specification.
\item[\textbf{P7}] (Prop.~\ref{prop:distortion}) Even where an assistant \emph{can}
condition on declared specification, terse users should still be asked less than
the assistant's own (possibly already-biased) per-type preference would predict,
specifically because of the incentive-compatibility constraint---a distortion
that should shrink as the population's cost heterogeneity ($\kappa_H-\kappa_L$)
shrinks, holding bias fixed (Remark on compounding).
\end{enumerate}

\section{Empirical Strategy (Summary)}

\textbf{Simulation.} Instantiate LLM-simulated users with heterogeneous cost
types $\kappa$ interacting with assistants implementing Model I's derived
$a^{SE}$, Model II's numerically-solved menu, and naive baselines (fixed
confidence threshold, ignoring both stakes and type). Verify
Proposition~\ref{prop:firstbest}'s threshold numerically, verify
Proposition~\ref{prop:aSE}'s closed form against a grid search over $a$, and
verify Proposition~\ref{prop:nodistortion}/\ref{prop:distortion} by solving the
constrained menu optimization directly and checking that the low type's bundle
sits at its unconstrained optimum while the high type's does not, with the gap
shrinking as bias is set to zero.

\textbf{Human study.} A between-subjects manipulation of ask policy (pooling vs.\
menu-conditioned; biased vs.\ unbiased friction weighting) across repeated coding
tasks, measuring specification length and, where feasible, self-reported
terseness/verbosity as a proxy for $\kappa$, to test P1$'$, P6, and P7 causally
rather than correlationally.

\section{Discussion and Limitations}
\label{sec:discussion}

The linear-risk approximation to the conjunctive success function is the main
tractability assumption carried through Sections~\ref{sec:stackelberg}--\ref{sec:compare};
we expect the qualitative results (threshold first best, no-distortion-without-bias,
downward distortion under bias) to survive the exact conjunctive form, since
they rely on convexity/monotonicity properties rather than the specific
functional form, but this should be checked numerically before the claims are
stated without qualification. The two-type restriction in Model II avoids
ironing/continuum-type complications at the cost of generality; a continuum
version using the standard Mirrlees technique is the natural next step and would
let $\kappa^*$ from the first-best benchmark interact directly with the screening
distortion. We also assume truthful revelation of specified attributes (no
cheap-talk vagueness within what is said) and a single leader (no competition
among assistants, which would itself be an interesting extension---competition
plausibly disciplines bias directly, another testable channel).

Two results in this model (the corner-solution character of unbiased pooling,
Corollary~\ref{cor:corner}; and the active constraint set in the biased
screening problem, the remark following Proposition~\ref{prop:distortion}) were
corrected after a companion numerical implementation solved the underlying
optimizations directly rather than trusting a hand-derived proof sketch by
analogy with classical mechanism design. Both corrections trace to the same
root cause---importing a technique that assumes the principal's payoff is
reduced by conceded surplus, into a setting where the principal's payoff
instead contains the agent's utility directly---which suggests this is a
recurring hazard specific to modeling an assistant whose objective is a
(possibly discounted) function of user welfare rather than a literal transfer,
and future extensions of this model should treat any imported result from the
transfer-based mechanism-design literature as a conjecture to be checked
numerically rather than a result to be cited directly.

\bibliographystyle{plainnat}
\begin{thebibliography}{9}

\bibitem[von Stackelberg(1934)]{stackelberg1934}
Heinrich von Stackelberg. \textit{Marktform und Gleichgewicht}. Springer, 1934.

\bibitem[Myerson(1981)]{myerson1981}
Roger B.\ Myerson. Optimal auction design. \textit{Mathematics of Operations
Research}, 6(1):58--73, 1981.

\bibitem[Baron and Myerson(1982)]{baron-myerson1982}
David P.\ Baron and Roger B.\ Myerson. Regulating a monopolist with unknown
costs. \textit{Econometrica}, 50(4):911--930, 1982.

\bibitem[Mussa and Rosen(1978)]{mussa-rosen1978}
Michael Mussa and Sherwin Rosen. Monopoly and product quality. \textit{Journal
of Economic Theory}, 18(2):301--317, 1978.

\bibitem[Laffont and Tirole(1986)]{laffont-tirole1986}
Jean-Jacques Laffont and Jean Tirole. Using cost observation to regulate firms.
\textit{Journal of Political Economy}, 94(3):614--641, 1986.

\bibitem[Crawford and Sobel(1982)]{crawford-sobel1982}
Vincent P.\ Crawford and Joel Sobel. Strategic information transmission.
\textit{Econometrica}, 50(6):1431--1451, 1982.

\bibitem[Townsend(1979)]{townsend1979}
Robert M.\ Townsend. Optimal contracts and competitive markets with costly state
verification. \textit{Journal of Economic Theory}, 21(2):265--293, 1979.

\bibitem[Grossman(1981)]{grossman1981}
Sanford J.\ Grossman. The informational role of warranties and private
disclosure about product quality. \textit{Journal of Law and Economics},
24(3):461--483, 1981.

\bibitem[Sims(2003)]{sims2003}
Christopher A.\ Sims. Implications of rational inattention. \textit{Journal of
Monetary Economics}, 50(3):665--690, 2003.

\bibitem[Topkis(1978)]{topkis1978}
Donald M.\ Topkis. Minimizing a submodular function on a lattice.
\textit{Operations Research}, 26(2):305--321, 1978.

\end{thebibliography}

\vspace{1em}
\noindent\textit{Note: bibliographic entries are given from memory to anchor
citations conceptually and should be checked against primary sources
(volume/issue/page numbers) before submission.}

\end{document}
